\begin{remark}[2.6.3]
\label{IV.2.6.3-fr}
En los casos (i) y (ii), las conclusiones de (2.6.2) siguen siendo válidas
cuando solo se supone que $g$ es \emph{cuasifielmente plano} (2.3.3) y \emph{cuasicompacto};
en efecto, en virtud de (2.1.4), (\textbf I, 3.5.2) y (\textbf I, 9.1.13.1), todavía es posible
reducirse al caso en que $Y=S$, $Y'=S'$; la conclusión se deduce entonces de
(2.3.12). En los casos (iii) y (iv), las conclusiones siguen siendo válidas
cuando solo se supone que $g$ es \emph{cuasifielmente plano}, siempre que se suponga
ya que $f$ es cuasicompacto; en efecto, solo se utilizan entonces (2.3.10) y el
hecho de que $g$ es sobreyectivo. Por último, si $g$ es fielmente plano y
localmente de presentación finita, o si $g$ es sobreyectivo y $S$ es discreto,
la conclusión de (2.6.2) es válida cuando $\boldsymbol{P}$ es la propiedad:
\begin{enumerate}
\item[(iii \emph{bis})] \emph{ser un homeomorfismo sobre el subespacio imagen};
\end{enumerate}
esto se deduce, en efecto, de la demostración dada en (2.6.2) y de la
Observación (2.4.11).
\end{remark}

\begin{corollary}[2.6.4]
\label{IV.2.6.4-fr}
Con las notaciones de \textup{(2.6.1)}, supongamos que el morfismo $g:S'\to S$ es
fielmente plano y cuasicompacto. Consideremos, para un morfismo, la propiedad
de ser:
\begin{enumerate}
\item[\textup{(i)}] universalmente abierto;
\item[\textup{(ii)}] universalmente cerrado;
\item[\textup{(iii)}] cuasicompacto y universalmente bicontinuo;
\item[\textup{(iv)}] un homeomorfismo universal;
\item[\textup{(v)}] cuasicompacto;
\item[\textup{(vi)}] cuasicompacto y dominante.
\end{enumerate}
Entonces, si $\boldsymbol{P}$ designa una de las propiedades anteriores, para que $f$ tenga
la propiedad $\boldsymbol{P}$, es necesario y suficiente que $f'$ la tenga.
\end{corollary}

\oldpage[IV-2]{29}
Las propiedades (v) y (vi) se incluyen solo a título recordatorio, pues son
consecuencias de (1.1.4), (1.1.6) y (2.3.7). En cuanto a las demás, la condición
es necesaria en virtud de (2.4.3). Recíprocamente, supongamos, por ejemplo, que
$f'$ es universalmente abierto, y sea $Y_1\to Y$ un morfismo arbitrario;
pongamos $X_1=X\times_Y Y_1$ y $f_1=f_{(Y_1)}$. Por otra parte, pongamos $Y'_1=Y_1\times_Y Y'$, $X'_1=X'\times_{Y'}Y'_1=X_1\times_{Y_1}Y'_1$, $f'_1=f'_{(Y'_1)}=f_{1(Y'_1)}:X'_1\to Y'_1$;
como $g'=g_{(Y')}:Y'_1\to Y'$ es fielmente plano y cuasicompacto
(2.2.13 y 1.1.2), y $f'$ es universalmente abierto, $f'_1$ es abierto, y de
(2.6.2) se deduce que $f_1$ es abierto; por tanto, $f$ es universalmente
abierto. El mismo razonamiento se aplica a los demás casos.

Se observará también aquí que se puede sustituir «~fielmente plano~» por
«~cuasifielmente plano~» y que, cuando $g$ es además localmente de presentación
finita, o cuando $g$ es sobreyectivo y $S$ es discreto, se puede sustituir la
propiedad (iii) por:
\begin{enumerate}
\item[(iii \emph{bis})] \emph{universalmente bicontinuo}.
\end{enumerate}

